% from ega4-18-fr.tex
\oldpage[IV-4]{150}

\begin{proposition}[18.8.17]
\label{IV.18.8.17}
Soit $A$ un anneau local noethérien. Si $A$ est universellement caténaire, il en est de même de ${}^{hs}\!A$.
\end{proposition}

La démonstration est la même que celle de (18.7.5).

\begin{proposition}[18.8.18]
\label{IV.18.8.18}
\begin{enumerate}
  \item[(i)] Toute limite inductive filtrante d'anneaux strictement locaux (où les homomorphismes de transition sont locaux) est un anneau strictement local.
  \item[(ii)] Le foncteur $A\leadsto{}^{hs}\!A$ dans la catégorie des anneaux locaux (où les morphismes sont les homomorphismes locaux) commute aux limites inductives filtrantes.
  \item[(iii)] Tout anneau strictement local $A$ est limite inductive d'un système inductif filtrant d'anneaux strictement locaux noethériens (les homomorphismes de transition étant locaux).
\end{enumerate}
\end{proposition}

Les démonstrations sont calquées sur celles de (18.6.14).

\subsection{Fibres formelles des anneaux noethériens henséliens.}
\EGAReaderTarget{ega4.c04.subsec.18-9}{subsection:IV.18.9}

\begin{theorem}[18.9.1]
\label{IV.18.9.1}
Soit $A$ un anneau local noethérien hensélien, dont les fibres formelles sont géométriquement normales. Alors les fibres formelles de $A$ sont géométriquement intègres (donc géométriquement connexes).
\end{theorem}

Les deux assertions de l'énoncé sont en fait équivalentes, un préschéma localement noethérien connexe et normal étant intègre. Comme toute $A$-algèbre intègre finie est un anneau local hensélien (18.5.9 et 18.5.10), il résulte de (7.3.16.2) que l'on est ramené à prouver que si l'on suppose de plus $A$ \emph{intègre}, alors la fibre du morphisme $\Spec(\widehat A)\to\Spec(A)$ au point générique de $\Spec(A)$ est \emph{intègre}, ce qui résulte du

\begin{corollary}[18.9.2]
\label{IV.18.9.2}
Soit $A$ un anneau local noethérien hensélien intègre, dont les fibres formelles sont géométriquement normales. Alors $\widehat A$ est intègre.
\end{corollary}

En effet, il résulte de (18.8.16) que $A$ est unibranche, donc le théorème résulte de l'hypothèse faite sur les fibres formelles de $A$ et de (7.6.3).

\begin{corollary}[18.9.3]
\label{IV.18.9.3}
Sous les hypothèses de (18.9.2), le corps des fractions $L$ de $\widehat A$ est une extension séparable du corps des fractions $K$ de $A$, et $K$ est algébriquement fermé dans $L$.
\end{corollary}

Cela résulte de ce que la fibre du morphisme $\Spec(\widehat A)\to\Spec(A)$ au point générique de $\Spec(A)$ est géométriquement intègre et de (4.3.2) et (4.3.5).

\begin{corollary}[18.9.4]
\label{IV.18.9.4}
Soit $A$ un anneau local noethérien, hensélien et dont les fibres formelles sont géométriquement normales (par exemple un anneau local noethérien hensélien et excellent), et posons $A'=\widehat A$. Soient $X$ un $A$-préschéma, $X'=X\otimes_A A'$, $g:X'\to X$ la projection canonique. Alors l'application $U\mapsto g^{-1}(U)$ est une bijection de l'ensemble des parties à la fois ouvertes et fermées de $X$ sur l'ensemble des parties à la fois ouvertes et fermées de $X'$.
\end{corollary}

Comme $g$ est un morphisme fidèlement plat et quasi-compact, on sait (2.3.12) que la topologie de $X$ est \emph{quotient} de celle de $X'$ par la relation d'équivalence définie par $g$. Tout revient donc à voir qu'une partie ouverte et fermée $U'$ de $X'$ est \emph{saturée} pour cette relation. Or, comme le morphisme $\Spec(A')\to\Spec(A)$ a ses fibres géométriquement connexes (18.9.1), les fibres de $g$ sont connexes, d'où aussitôt notre assertion.

En particulier, si $X$ est localement connexe (ce qui sera le cas si $X$ est localement noethérien), alors, pour toute composante connexe $U$ de $X$ (qui est ouverte et fermée

